\documentclass[../main-detail.tex]{subfiles}
\begin{document}
\chapter{概要}
\paragraph{本書に関する説明}
本書はコノメノの仕様書である．コノメノの解説書は，以下の3冊構成とする予定である．
\begin{itemize}
    \item \emph{コノメノ詳説}(開発者・上級者用)：詳細な理論から実践までを網羅する仕様書．
    \item \emph{コノメノ素描}(初心者用)：分かりやすく素朴に文法事項を整理した入門書．
    \item \emph{コノメノ入門}(初心者用)：物語形式で本文を読み進める入門書．
\end{itemize}
\begin{remark}
    旧来は『コノメノ入門』と『コノメノ詳説』の2冊構成である．『コノメノ入門』は本文用の小説執筆のためのメタフィクション理論待ちで停滞しており，『コノメノ素描』はそのバイパスである．
\end{remark}

\section{モチベーションとアプローチ}
\paragraph{P-side \& S-side} 言語学のサブフィールドを
\begin{itemize*}
    \item[] P-side：phonetics \& phonology，
    \item[] S-side：syntax \& semantics，
\end{itemize*}
と分けて呼ぶ慣例にならって，コノメノの制作部門もこのように分類する．P-sideは芸術言語志向，S-sideは工学言語志向が強い．両者は独立に開発できる：音韻は文法（意味論・オントロジー）の設計に関与しない．
制作部門は以下の通り：
\begin{itemize}
    \item \emph{P-side}：
    \begin{itemize*}
        \item \emph{音韻部門}
        \item \emph{文字部門}
    \end{itemize*}
    \item \emph{S-side}：
    \begin{itemize*}
        \item \emph{意味論部門}
        \item \emph{オントロジー部門}
    \end{itemize*}
\end{itemize}
\memo{
文字は音韻から独立に発展しうるため，文字部門をP-sideに置くこの分類は暫定である．
}
\memo{
P-side, S-sideという用語をそのまま使うべきかは要検討．意味のずれとスラング感．
}

\paragraph{P-side -- 音韻部門} 音の印象（音象徴・共感覚など）に基づく造語を中心に，語彙全体に固有の表情があり，日常的に使ってもその表情が保たれることを目指す．音韻と造語の章では，現行体系と制作の目標・方法を区別し，制作時と使用・再評価時の不満を体系の改訂へ戻す．統計は偏りを探す資料として用い，既存の分布や自然な発音に合うことだけを完成目標にはしない．

\paragraph{文字部門}
% @issue KONO-0039
\todo{あとで書く}

\paragraph{S-side -- オントロジー部門} 「日常的な推論を形式的に表現すること」を目標とする．
\suppl{言語名の由来も大まかにこれに基づいている：``konomeno''という名称は，``kono''（概念）と``meno''（法則）からなり，文字通りには，コノメノとは「概念の法則を記述する言語」である．}
ここでの「日常的な推論」は語の意味に密着した推論を指し，「雨が降ると地面が濡れる」のような経験的知識を要するいわゆる常識推論よりも弱い．
\memo{
長い．ここで話すべき話ではなくないか．あと主題と別に書くべきでもない．編成をするためのプロンプトが欲しい．美意識がない．あるいは文章の書き方の問題でこちらの意図の深い部分まで伝わっていないか．
}
このためには，
\begin{enumerate*}
    \item （意味論部門）自然言語の（形式）意味論，
    \item （オントロジー部門）知識工学をはじめとした旧来の人工知能研究，および哲学における形式存在論，
\end{enumerate*}
がよくモチベーションに合致する．意味論部門は言語表現を，オントロジー部門は概念体系をそれぞれモデル化する．

Montague意味論をはじめとする古典的なモデル論的意味論は，概念の中身に深入りしない．その枠内でできるのは
\begin{align*}
    \infer{\psi}{\phi\quad \phi \to \psi} \qquad \infer{\exists x\phi(x)}{\phi(a)}
\end{align*}
のような純粋に論理的な推論に限られる\footnote{この種の推論の妥当性すら，古典論理の「ならば」が関連性などの面で自然言語の「ならば」と乖離しているという指摘のもとでは断言できない．}．しかし日常的な推論の多くは概念の性質に依存する：認識論理の公理4
\begin{exe}
    \ex 「$\phi$を信じている」$\vdash$「『$\phi$を信じている』を信じている」
\end{exe}
を立てるには「信じる」という概念の性質を，
\begin{exe}
    \ex 「金庫に財布が入っている」「財布にお金が入っている」$\vdash$「金庫にお金が入っている」
\end{exe}
を認めるには「入っている」の推移性を知らなければならない．こうした概念の性質（$=$概念体系の百科事典的知識）を語彙と公理として際限なく整備していくのがオントロジー部門であり，表現の豊かさは造語だけでなくこのモデル化の蓄積にも比例する．

モデル化の根拠は，できるだけ言語の観察（日常的な推論の観察）から得る．例えば
\begin{exe}
    \ex 「$x$は楽しい」$\equiv$「$x$は楽しさがある」
    \ex 「$x$は大きい」$\not\equiv$「$x$は大きさがある」
\end{exe}
という対比（「大きい」は基準超過を意味するが「大きさがある」はそうでない）は，同型に見える形容詞の意味的な差を示す．また，モデル化は一意でない．概念の近似\footnote{「$x$は本来ロール概念だが，簡単のためにコンテキストは無視してプレーヤーのみで考える」といった近似．}によって適用範囲の異なる複数のモデルが得られ，この非一意性は形式意味論章で超内包性の表現手段として利用する（\ref{...}節）．

% @issue KONO-0005
\todo{オントロジーの定義．辞書の構造の話．}
\memo{S-sideの主題のToaq比較（オントロジーを持たない）がこの節の存在に依存している．優先度高．}

\paragraph{意味論部門}
形式意味論には「言語表現がどのような形式的対象（論理式など）に翻訳されるか」と「統語構造と意味構造がいかに対応するか」という2つの関心がある．コノメノは前者に重点を置く．文法は基本的に意味構造に忠実であるように設計されているため論理式らしい見た目をしており，ふりかけ程度に，自然言語らしさや日常会話で使う文を簡潔に表現できる工夫がされている．この工夫は実際の言語における「統語構造と意味構造の対応」を真似たものではない．むしろ，誤った一般化から生じる奇妙さを面白がって``誤り''を取り入れている．
\memo{オントロジー部門に例が3つあるのに対しこの段落には例がなく，アンバランス．文語の文例を1つ足すか検討（形式文法章と重複するので保留）．}

コノメノには，\textbf{文語コノメノ}（胡：keltokono）と\textbf{口語コノメノ}（胡：kuyskono）がある．文語は形式言語であり，省略なしで書けば構文的曖昧性がない．いくつかの体系があるが，基本的に表記は記号的であり，発話するための音韻はない．口語は日常会話で使うための言語である．詳細は形式文法章（\ref{...}）で扱う．

\paragraph{S-sideの主題}
S-sideの制作主題は次の2つである．
\begin{enumerate}
    \item \emph{人工言語の観点から}：形式意味論と形式オントロジーを本格的に導入した人工言語を実現すること．
    \item \emph{意味論・オントロジーの観点から}：複数の意味論・オントロジーを1つの言語に組み込めることを実証すること，すなわちそのためのインターフェイスを設計すること．
\end{enumerate}
主題2の動機は2つある．第一に\emph{大規模性}：1つの理論の守備範囲は限定されている——形式意味論は伝統的に，時制・複数性・態度動詞といった現象ごとに断片（fragment）を切り出し，個別の理想化のもとで理論を立ててきた——ため，言語全体を覆うには複数の理論の合成が必要になる．第二に\emph{非一意性}：同じ現象に対して粒度や近似の異なる複数のモデルが正当に共存する（オントロジー部門で述べた通り，これは超内包性の表現手段でもある）．この2方向に対応してインターフェイスも2つ実装されている：水平方向（現象ごとの理論の合成）はモジュール構成（体系P＋線形文・指示詞・談話モジュール）と引き下げ対応（\ref{...}）が，垂直方向（粒度の共存）は冠詞$M$（表現）・$C$（状況）・$K$（命題）と段の間の粗視化写像（\ref{...}）が担う．なお，理論の合成は並べるだけでは済まない——枠組み同士は型や存在論的コミットメントで衝突する——ため，衝突の調停までがインターフェイスの仕事である．

先行する試みは，意味論とオントロジーのどちらか一方に留まっている．論理的人工言語の系譜では，Lojbanは機械検証された文法でパースの一意性を保証するが，パース木を論理式へ写す公式の手続き（つまり意味論）を持たない\cite{cll}．Toaqは型付き$\lambda$計算・イベント意味論・可能世界による標準的な合成的意味論を言語全体に通した\cite{toaq-refgram}が，意味論的在庫は個体・イベント・世界・時刻に留まり，概念分類や百科事典的公理（オントロジー）を持たない．逆に，語彙を分類体系から構成する言語（Wilkinsの哲学的言語\cite{wilkins1668}，Ro）は分類だけを持ち，形式的な意味論を持たない．

両方を備える学術側の体系ではOntological Semantics\cite{nirenburg-raskin2004}が近い．そのMikrokosmos実装は，約6000概念の言語中立オントロジー，言語別レキシコン，テキスト意味表現（TMR），推論処理を組み合わせる．ただし，自然言語を解析・生成するための内部体系であり，オントロジーを人が使う言語の語彙とするコノメノとは異なる．UNLとCycも同じく機械用の知識表現であって，人が話す言語ではない．

人工言語全体として最も近いのはIEMLである\cite{levy2010ieml,levy2023ieml}．IEMLは約3000語の辞書と規則文法を持ち，記号列の代数的構造に語の範列関係を符号化して概念グラフを計算し，複数のオントロジーの相互運用を目指す．ただし，その「計算可能な意味論」は主に言語内の概念関係の計算であり，真理条件や指示を与えるモデル論は別の層に置く．また，読み書きする意味コードとして設計され，話し言葉ではない．コノメノはこれに対し，人が使う言語にモデル論的意味論とオントロジーを組み込み，複数のモデルと粒度を同じ文法から扱う．

主題2のインターフェイスという発想自体にも先行例がある．計算意味論のMRSは，自らを意味理論ではなく複数の意味表現を接続するメタレベルの枠組みと規定する\cite{copestake2005mrs}し，オントロジー工学のDOLは，異なる論理で書かれたオントロジー群を関係づけるためのメタ言語である\cite{mossakowski2015dol}．ただし，人工言語そのものをこのインターフェイスとして設計した例は管見の限りない．

\paragraph{立ち位置}
コノメノは\emph{理論創作}——空想的な現象について数理モデルを立てる（しばしば現象-理論のペアを同時に設計する）活動——の作品である．したがって新規性の主張は第一に創作作品としてのものであり，個々の理論的道具はほとんどが既存研究からの輸入である．学術的な位置づけは，何を輸入し，どの組み合わせが先行例を持たないかという上記のような事実の報告に留める．


% @issue KONO-0006
\todo{文語と口語の話もっと書く．これは形式文法の体系が確定してから．LojbanやIthkuilとの比較も書く．}

\section{記法の凡例}
本書で使用する記法について，いくつかの凡例を示す．特に断りがない限り，以下の記法を用いる．
\subsection{アルゴリズム関連の記法}
本書では，複数の章で文字列の操作や擬似コードを用いるため，文字列・リストの基本記法をここで統一しておく．これらの記法はPython風の見た目を採用しつつ，数式として読めるように整理したものである．

\begin{definition}[文字列とリストの基本記法]
以下の記法を用いる．
\begin{enumerate}[label=(\arabic*)]
    % \item アルファベットを$\Sigma$とし，文字列全体を$\Sigma^*$とする．文字列は$s,t,u\in\Sigma^*$のように表す．
    \item リストは$X,Y,Z$のように表し，要素列$[x_0,\ldots,x_{n-1}]$で書く．空リストは$[]$，空文字列は$\varepsilon$とする．
    \item 長さは，文字列・リストともに$|\bullet|$で表す．
    \item 文字もしくは文字列の集合$\Sigma$に対して，Kleene 閉包を$\Sigma^* = \bigcup_{n\ge 0} \Sigma^n$とする．
    \item 連結は$\circ$で表す（例：$s\circ t,\; X\circ Y$）．意味が明確であれば$\circ$を省略して書いても良い．
    \item リストあるいは文字列の集合$X, Y$について，$X \circ Y = XY = \{x\circ y \mid x\in X, y\in Y\}$とする．
    \item 要素アクセスは$X_i$もしくは$X[i]$（$0\le i<|X|$）とする．特に，\textbf{0-indexed}であることに注意．文字列も同様に$s_i$で$i$文字目を表す．$i < 0$の場合は後ろからアクセスするものとする．
    \item スライスは半開区間で$X_{i:j}$もしくは$X[i:j]$と書き，$[X_i,\ldots,X_{j-1}]$を表す（$0\le i\le j\le |X|$）．文字列$s_{i:j},\; s[i:j]$も同様．
    \item 集合的な記法として，$x\in X$で「$x$がリスト$X$に含まれる（少なくとも1回現れる）」を表す．必要に応じて，集合化を$\mathrm{set}(X)=\{x\mid x\in X\}$と書く．
    \item リスト内包 $[f(x)\mid x\in X,\; P(x)]$は，$X$を左から走査し，条件$P(x)$を満たす要素だけに$f$を適用して得られるリストを表す．
    \item 段階内包 $[k\mid x\in X,\; P(x)\mid k\in x,\; Q(x,k)]$は，まず$x\in X$を走査かつ$P(x)$を満たす要素を選び，次に各$x$の内部の$k\in x$を走査かつ$Q(x,k)$を満たす要素を選んで得た$k$を順に並べる記法とする．
\end{enumerate}
\end{definition}
\begin{remark}
$\in$は「集合の所属」と「リスト走査対象」の両方に使うが，文脈で区別できるものとする．
\end{remark}
\begin{definition}[その他のアルゴリズムでよく使う操作]
    以下の記法を用いる．
    \begin{enumerate}[label=(\arabic*)]
        \item 三項演算子$x\;\text{if}\; P\;\text{else}\; y$は，条件$P$が真のとき$x$を，偽のとき$y$を返す．
    \end{enumerate}
\end{definition}

\subsection{意味論関連の記法}
\begin{definition}[論理式の記法]
以下の記法を用いる．
\begin{enumerate}[label=(\arabic*)]
    \item 意味論の文脈で特に断りがない限り，$L$は1階述語論理の言語を表す．
    \item $L$-構造は$\mathfrak M,\mathfrak N$のように表す．$L$-構造全体(真クラス)は$\text{$L$-Strc}$と書く．
    \item $L$-論理式は$\phi,\psi$のように表す．$L$-論理式全体は$\text{$L$-Fml}$と書く．$L$-閉論理式 (文) は$\text{$L$-CFml}$と書く．
    \item $L$-変数は$x,y,z$のように表す．$L$-変数全体は$\text{$L$-Vars}$と書く．
    \item 充足関係は$\vDash$で表す．充足関係の左辺にはいくつかのパターンがある．
    \begin{enumerate}
        \item $\mathfrak M \vDash \phi$：構造$\mathfrak M$が閉論理式$\phi$を充足する．
        \item $\mathfrak M, \sigma \vDash \phi$：$\phi$は自由変数を含むかもしれない論理式であり，$\sigma$は変数割り当てである．$\sigma$全体は$\text{$L$-Vars}$から$|\mathfrak M|$への有限部分関数である．$\mathfrak M$は$\phi$を$\sigma$のもとで充足する．
        \item $M, \rho, \sigma \vDash \phi$：構造$\mathfrak M$は，台集合$|\mathfrak M|$と述語記号・関数記号の割り当て$\rho$の組と同一視できる($\rho$を語彙割り当てとも言う)．
        \item $\chi \vDash \phi$：前提$\chi$が結論$\phi$を導く．完全性定理のもとで$\chi \vDash \phi$は帰結関係$\chi \vdash \phi$と同値である．
    \end{enumerate}
    \item 帰結関係は$\vdash$で表す．$\chi \vdash \phi \iff 
        \forall \mathfrak M\forall\sigma.\; \mathfrak M,\sigma \vDash \chi \to \mathfrak M,\sigma \vDash \phi$である．
\end{enumerate}
\end{definition}
\begin{definition}[関数の記法]
    以下の記法を用いる．
    \begin{enumerate}[label=(\arabic*)]
    \item 集合$X$から集合$Y$への関数全体を$Y^X$または$\text{Func}(X,Y)$または$\mb{X,Y}$のいずれかで表す．
    \item 集合$X$から集合$Y$への部分関数の集合を$\text{Pat}(X,Y)$と書く．
    \item 部分関数$f\in\text{Pat}(X,Y)$に対して，$\dom(f) = \{x\in X \mid \exists y\in Y, f(x)=y\}$を$f$の定義域とする．
    \item 集合$X$から集合$Y$への有限部分関数の集合を$\text{Fin}(X,Y) = \{f\in\text{Pat}(X,Y) \mid |\dom(f)| < \infty\}$とする．
    \item 部分関数$f\in\text{Pat}(X,Y)$と$A \subseteq X$に対して，$f|_A = \{\langle x,y\rangle \in f \mid x\in A\}$とする．
    \item 部分関数$f\in\text{Pat}(A,B)$と$g\in\text{Pat}(C, D)$に対して，
    \begin{enumerate}
        \item $f \vlr g = f \cup g|_{C \setminus \dom(f)}$
        \item $f \vll g = f|_{A \setminus \dom(g)} \cup g$
    \end{enumerate}
    とする．$f \vlr g$は$f$の$g$による拡張，$f \vll g$は$f$の$g$による更新と解釈できる．
\end{enumerate}
\end{definition}
\subsection{オントロジー関連の記法}
\todo{後で書く}
\begin{enumerate}
    \item 個体，実例，インスタンス，トークンを等価な用語として扱う．主に個体を使う．
    \item 種，クラス，タイプ(，普遍者？)を等価な用語として扱う．
    \item スロット，フィールド，属性を投下な用語として扱う．
    \item クラス間の包含を$\sqsubseteq$で表す．
    \item 
    \item $\cdots$
\end{enumerate}
\begin{figure}[H]
    \centering
    \begin{forest}
        dir tree
        [$a_0$
          [$a_{00}$
            [$a_{000}$]
            [$a_{001}$,name=001]
          ]
          [$a_{01}$
            [$a_{010}$]
            [,name=011,dashed]
          ]
        ]
        \draw [dashed] (001.east) to[out=0,in=0] (011.east);
    \end{forest}
\end{figure}
\subsection{音韻関連の記法}
\todo{後で書く}
\subsection{用語}
\todo{
日本語訳がそれほど一般的ではないかもしれない用語
\begin{enumerate}
    \item 原子：atom, urelement
    \item 署名：signature
    \item 超内包性：hyperintensionality
\end{enumerate}
}
\section{アプリオリ性について}
% cf.セレニズム
実務的に，国の呼び方や単位系を同語するのにあたって特定の文化への依存の仕方について考える必要が出てくる．文化中立的な記述は難しい．例えば単位系について，SI単位系を「SI単位系モジュール」として定義することは可能だが，国の呼び方について同じことをすることが難しい．

\begin{enumerate}
    \item 借用に対する立場
    \item 文字に対する立場(なぜラテン文字綴りを使うのか，など)
    \item 
\end{enumerate}


\todo{あとで書く}

\section{今後の展望}
\memo{展望とはいえ，曖昧で薄い内容は書きたくない．}
\begin{enumerate}
    \item プログラミング言語にする．コノメノは意味論が目的に対して過度に煩雑に思えるから，他の言語を設計し直して行う可能性が高い．
    \item メタフィクションに関する理論の制作．本言語がどの文化に依拠して作られているのか，OCにとってどのような立ち位置にあるのかを明確にすること．
    \item 架空世界の物理体系を制作し言語と接続すること．命令により行為がなされること，乱数がプログラムに副作用を与えることをモデル化する．
    \item 文法・音韻における変遷のモデル化．
\end{enumerate}
\end{document}
